% !TeX program = xelatex
% ============================================================
% pp/pp.tex — Project Proposal (IEEE conference layout, English)
%
% Требования ОП ПИ (Project Proposal 2026):
%   • объём 1500–3000 слов (без abstract, keywords и references);
%   • Abstract 100–150 слов (в объём не входит);
%   • Keywords: 5–10, В АЛФАВИТНОМ ПОРЯДКЕ;
%   • Introduction ≈300 слов; Main Body 1200–1500 слов
%     (Literature Review, Methodology, Results ≈100 слов);
%   • Conclusion ≈200 слов;
%   • Библиография: МИНИМУМ 10 источников, из них ≥8 на английском
%     (до 2 неанглоязычных допускается), стиль IEEE ([1], [2], …);
%   • антиплагиат: не более 20% заимствований;
%   • только английский язык, две колонки, Times New Roman 10 pt
%     (поля 1,5 см слева/справа и 2 см сверху/снизу даёт класс IEEE);
%   • научный (академический) стиль: пассив и безличные конструкции,
%     без первого лица (I/we) и разговорной/«простой» лексики;
%   • текст пишется САМОСТОЯТЕЛЬНО: автоматический перевод всего PP
%     генеративными моделями или переводчиками ЗАПРЕЩЁН регламентом;
%   • если использовался генеративный ИИ — его нужно ЗАДЕКЛАРИРОВАТЬ
%     (см. блок REFERENCES ниже): сноска в тексте + запись в списке.
%
% Жёлтые блоки \hseFill{…} — места, которые нужно заменить своим
% текстом. Данные автора приходят из настроек проекта через макросы
% \hse… (common/meta.tex перегенерируется перед каждой сборкой).
% ============================================================
\documentclass[conference]{IEEEtran}
\IEEEoverridecommandlockouts

% ── Пакеты ──────────────────────────────────────────────────
\usepackage{cite}
\usepackage{amsmath,amssymb,amsfonts}
\usepackage{graphicx}
\usepackage{xcolor}
\usepackage{url}
\usepackage{hyperref}
\usepackage[english]{babel}
\usepackage{balance}  % выравнивание колонок на последней странице

% Шрифты — из настроек проекта (Настройки → Шрифты), с фолбэком на
% свободные клоны: сборка не падает, даже если шрифта нет в контейнере.
\newcommand{\hseDefaultMono}{Courier New}
\input{../common/fonts}

% Данные проекта (автогенерируемые макросы \hse… и \hseFill)
\input{../common/meta}

% Кликабельные ссылки
\hypersetup{
    unicode=true,
    breaklinks=true,
    bookmarksopen=true,
    bookmarksnumbered=true,
    colorlinks=true,
    linkcolor=black,
    urlcolor=blue,
    citecolor=black
}

% Перенос длинных URL в библиографии
\def\UrlBreaks{\do\/\do-\do.\do=\do?\do\&}
\sloppy

\begin{document}

\title{\hseTitle}

\author{\IEEEauthorblockN{\hseAuthorNameEn}
\IEEEauthorblockA{\textit{\hseFacultyEn} \\
\textit{\hseUniversityEn}\\
\hseCityEn \\
\hseAuthorEmail}}

\maketitle

% ── Abstract: 100–150 слов, что сделано/планируется — кратко ──
\begin{abstract}
\hseFill{Briefly state the problem you address, why it matters, what
you propose to build, how you will evaluate it, and what results you
expect. Keep it 100 to 150 words and write it after the rest of the
proposal is done.}
\end{abstract}

% ── Keywords: 5–10 шт. (жёсткий минимум методички — 3), строго по
%    алфавиту, через запятую. Первое слово с заглавной, остальные со
%    строчной — кроме аббревиатур и имён собственных (IoT, Kubernetes).
\begin{IEEEkeywords}
\hseFill{five to ten keywords, in alphabetical order, comma separated}
\end{IEEEkeywords}

% ============================================================
% INTRODUCTION — ≈300 слов. Ответьте по порядку на вопросы:
%   1) какова область исследования (1–3 предложения);
%   2) какую проблему решаете и почему она важна (2–4 предложения);
%   3) как планируете решать (1–2 предложения);
%   3a) сформулируйте ЦЕЛЬ (и задачи) и/или исследовательские вопросы;
%   4) на какие работы опираетесь;
%   5) в чём новизна подхода по сравнению с существующими;
%   6) какие результаты ожидаются (1–4 предложения);
%   7) как организована статья (3–4 предложения).
% Проектная (не исследовательская) ВКР: вместо Literature Review —
% анализ стейкхолдеров/целевой аудитории и обзор существующих
% продуктов; вместо Methodology — раздел «Methods and/or Product
% Design». Остальные разделы — общие.
% ============================================================
\section{Introduction}

The research area of this work is \hseFill{name the research area in
one to three sentences}.

\hseFill{State the problem you are trying to solve and explain why it
is important. Support the claim with evidence from the literature}
\cite{samplebook}.

This work proposes \hseFill{outline your solution approach in one or
two sentences}. The approach builds upon \hseFill{name the studies or
practices your research is based on} \cite{sampleweb}. The key novelty
lies in \hseFill{explain what makes your approach new compared to the
existing ones}.

The aim of this work is \hseFill{state the aim of the project and the
objectives (or research questions) through which it is achieved}.

Expected results include \hseFill{describe the expected results in one
to four sentences}.

The rest of this proposal is organized as follows: Section II reviews
the existing solutions, Section III describes the methodology, Section
IV outlines the expected results, and Section V concludes the paper.

% ============================================================
% MAIN BODY — 1200–1500 слов суммарно (разделы II–IV).
% Пишите СВЯЗНО: используйте linking words (However, Moreover,
% Therefore, In contrast, As a result…) и делите текст на короткие
% абзацы — логичность, связность и абзацное членение оценивают.
%
% Literature Review: введение в аспект → что говорит литература →
% мини-вывод, чем ваша работа продвигает известное; затем следующий
% аспект. Текст должен быть АНАЛИТИЧЕСКИМ, а не описательным: не
% пересказ источников, а их сопоставление; поясняйте, какие источники
% использованы и как. Сравнительная таблица аналогов приветствуется.
% Проектный трек: здесь — анализ стейкхолдеров/целевой аудитории и
% обзор существующих продуктов-аналогов.
% ============================================================
\section{Literature Review}

\hseFill{Introduce the first aspect of your research area, summarize
what the literature says about it, and add a mini-summary explaining
how your work advances what is already known. Then move on to the next
aspect. Cite every source you mention}
\cite{samplebook, sampleweb, samplejournal, sampleimage}.

\hseFill{Compare the closest existing solutions against your proposal.
A comparison table is a compact way to do it; reference it from the
text}, as shown in Table~\ref{table:comparison}.

\begin{table}[h]
\centering
\caption{Comparative analysis of existing solutions}
\label{table:comparison}
\small
\begin{tabular}{lccc}
\hline
\textbf{Criterion} & \textbf{Tool A} & \textbf{Tool B} & \textbf{Proposed} \\
\hline
Criterion 1 & + & - & + \\
Criterion 2 & - & + & + \\
Criterion 3 & - & - & + \\
\hline
\end{tabular}
\end{table}

% ============================================================
% Methodology: что изучаете, какие гипотезы проверяете, как устроен
% эксперимент, какие переменные измеряете и почему, какие ПО, ДАННЫЕ
% и оборудование используете и ОТКУДА они; как и ПОЧЕМУ выбраны методы
% анализа и выборка; ссылки на литературу; ожидаемые трудности.
% Проектный трек: раздел называется «Methods and/or Product Design».
% ============================================================
\section{Methodology}

\hseFill{Describe what you study and what hypotheses you test. Explain
how the experiment or system is designed and what assumptions are
made.}

\hseFill{List the technology stack, the data and the equipment or
infrastructure you use and where they come from, and justify the
choice of analysis methods and sampling} \cite{samplesoft}.

\hseFill{Mention the difficulties you expect to face and how you plan
to mitigate them.}

% ============================================================
% Results — ≈100 слов: главные ожидаемые результаты, второстепенные,
% подтверждающие данные; результаты, противоречащие гипотезе, тоже
% упоминаются.
% ============================================================
\section{Expected Results}

\hseFill{Summarize the main expected results and deliverables of the
project in about one hundred words. Mention the metrics you will use
to validate them and what would falsify your hypothesis. Make sure
each result traces back to an objective from the Introduction and a
method from the Methodology.}

% ============================================================
% CONCLUSION — ≈200 слов: кратко повторить главные положения, их
% значимость, возможные улучшения и рекомендации для будущей работы.
% ============================================================
\section{Conclusion}

\hseFill{Briefly revisit the most important points of the proposal and
the significance of the expected findings. Relate the outcomes to the
aim and objectives stated in the Introduction. Suggest improvements
and recommendations for future work. Aim at about two hundred words.}

% Подсчёт слов: тулбар «Превью» показывает число слов после каждой
% сборки (texcount) — впишите его сюда перед сдачей.
\vspace{2em}
\begin{center}
\rule{0.5\linewidth}{0.5pt}\\[0.5em]
\textit{Word Count \hseFill{fill in before submission}}\\
\textit{(excluding title, abstract, keywords, references, appendices,
and the AI-use declaration)}\\
\rule{0.5\linewidth}{0.5pt}
\end{center}

% ============================================================
% REFERENCES — МИНИМУМ 10 источников, из них ≥8 на английском
% (до 2 неанглоязычных допускается), стиль IEEE.
%
% Правила IEEE:
%   • источники нумеруются В ПОРЯДКЕ ПЕРВОГО ЦИТИРОВАНИЯ в тексте,
%     НЕ по алфавиту; повторная ссылка использует тот же номер;
%   • в список включайте ТОЛЬКО то, на что есть \cite в тексте;
%   • ресурсы Web 2.0 (вики, форумы, блоги, соцсети) как источники
%     ЗАПРЕЩЕНЫ.
%
% Декларация ИИ (обязательна, если применялся генеративный ИИ):
%   • сноска в тексте — модель и промпт (в объём слов не входит),
%     например:  \footnote{Generated with ChatGPT (GPT-4o); prompt:
%     ``...''.};
%   • И отдельный пункт этого списка: название чата/модели, издатель,
%     промпт, дата обращения — например «ChatGPT (GPT-4o), OpenAI.
%     Prompt: ``...''. Accessed: ...».
%
% Ниже — образцы всех типов из методички (книга, статья в журнале,
% веб-страница, изображение, ПО). Замените их своими источниками и
% добавьте остальные по тем же образцам.
% ============================================================
\balance
\begin{thebibliography}{99}

% Книга: Автор(ы), Название, издание, издательство, год.
\bibitem{samplebook}
T. Brown, \textit{Introduction to Electronics}, 3rd ed., Pearson, 2020.

% Веб-страница: Автор(ы), «Заголовок,» Сайт, дата.
% [Online]. Available: URL. [Accessed: дата].
\bibitem{sampleweb}
K. Jacob, ``Future of climate change,'' Climate Talk, Mar. 20, 2024.
[Online]. Available: \url{https://techtalk.com/climate-future}.
[Accessed: Apr. 5, 2025].

% Статья в журнале: Автор(ы), «Название,» Журнал (сокр.),
% т., №, сс. xx--xx, месяц год.
\bibitem{samplejournal}
L. Chen and R. Patel, ``Advances in wireless technology,''
\textit{IEEE Trans. Commun.}, vol.~68, no.~4, pp.~123--130, Apr. 2021.

% Изображение: Автор(ы), «Название,» Источник, дата.
% [Online]. Available: URL. [Accessed: дата].
\bibitem{sampleimage}
M. Clark, ``Neural network diagram,'' Neuralgraphics, Feb. 10, 2023.
[Online]. Available: \url{https://neuralgraphics.com/nn-diagram}.
[Accessed: Apr. 5, 2025].

% Программное обеспечение: Название. (версия/год), издатель.
% Accessed: дата. [Online]. Available: URL.
\bibitem{samplesoft}
Antenna Products. (2011), Antcom. Accessed: Feb. 12, 2014. [Online].
Available: \url{http://www.antcom.com/documents/catalogs/L1L2GPSAntennas.pdf}

\end{thebibliography}

% ============================================================
% APPENDICES (факультативно) — вспомогательный/справочный материал
% (таблицы, графики, формулы), не критичный для основного текста; в
% подсчёт слов НЕ входит. Раскомментируйте при необходимости:
% \appendix
% \section{Supplementary material}
% \hseOptional{additional tables, figures or derivations}
% ============================================================

\end{document}
